% Backup material retained for discussion after the 16-slide spoken path.

\begin{frame}{Backup: Generation Engines Do Not Automatically Win}
  \small
  \begin{tabularx}{\textwidth}{llYYY}
    \toprule
    \textbf{Engine} & \textbf{Variant} & \textbf{MB/s} & \textbf{Ratio} & \textbf{Speedup} \\
    \midrule
    Transformers & + KV cache & \metric{0.0344} & 7.0492 & $1.00\times$ \\
    Transformers & No KV cache & 0.0030 & 7.0502 & $0.09\times$ \\
    vLLM & Default & 0.0130 & 7.0304 & $0.37\times$ \\
    TensorRT-LLM & Batch 256 & 0.0081 & 7.0499 & $0.023\times$ \\
    TensorRT-LLM & Batch 512 & 0.0065 & 7.9370 & $0.019\times$ \\
    \bottomrule
  \end{tabularx}

  % \vspace{0.8em}
  % \begin{block}{Takeaway}
  %   vLLM and TensorRT-LLM are excellent generation systems, but this workload exposes scheduling, synchronization, and score-capture costs that ordinary serving can hide \parencite{kwonEfficientMemoryManagement2023,Kipf2026WaitingTD}.
  % \end{block}
\end{frame}

\begin{frame}{Backup: Making coding Asynchronous}
  \centering
  \resizebox{0.98\textwidth}{!}{%
  \begin{tikzpicture}[
    box/.style={draw,rounded corners=2pt,minimum width=2.65cm,minimum height=1.05cm,align=center,font=\small},
    arr/.style={-{Stealth[length=5pt]},thick},
  ]
    \node[box,fill=blue!10] (vec) at (0,2.1) {\textbf{vectorize}\\$[B,|V'|] \rightarrow [B,3]$\\target intervals};
    \node[box,fill=mLightBrown!15] (compile) at (3.45,2.1) {\textbf{compile}\\Numba integer\\range updates};
    \node[box,fill=mDarkTeal!12] (streams) at (6.9,2.1) {\textbf{multistream}\\one $(low,high)$\\state per row};
    \node[box,fill=green!12] (packed) at (10.35,2.1) {\textbf{pack}\\write bytes directly,\\not Python bit lists};
    \draw[arr] (vec) -- (compile);
    \draw[arr] (compile) -- (streams);
    \draw[arr] (streams) -- (packed);

    \node[anchor=east,font=\scriptsize\bfseries] at (-0.35,0.35) {GPU};
    \node[anchor=east,font=\scriptsize\bfseries] at (-0.35,-0.55) {CPU};
    \foreach \i/\x in {0/0,1/3.0,2/6.0,3/9.0} {
      \node[draw,fill=blue!14,minimum width=2.65cm,minimum height=0.58cm,font=\scriptsize] at (\x+1.3,0.35) {infer step $t+\i$};
    }
    \node[draw=none,minimum width=2.65cm,minimum height=0.58cm,font=\scriptsize,text=black!45] at (1.3,-0.55) {idle};
    \foreach \i/\x in {0/3.0,1/6.0,2/9.0} {
      \node[draw,fill=mDarkTeal!14,minimum width=2.65cm,minimum height=0.58cm,font=\scriptsize] at (\x+1.3,-0.55) {encode step $t+\i$};
    }
    \draw[decorate,decoration={brace,amplitude=4pt},thick,mDarkTeal] (3.0,-1.05) -- (11.65,-1.05);
    \node[font=\scriptsize,mDarkTeal] at (7.3,-1.42) {background thread overlaps CPU coding with the next GPU call};
  \end{tikzpicture}
  }

  \vspace{0.35em}
  \small Pipelining is safe because compression already knows the next file token: inference for step $t+1$ does not wait for step $t$'s bits to be emitted.
\end{frame}

\begin{frame}{Backup: Encoder and Decoder Have Different Information}
  \begin{columns}[T,onlytextwidth]
    \begin{column}{0.47\textwidth}
      \begin{block}{Compression: $x_{t+1}$ is known}
        \begin{enumerate}
          \item Compute and record honest scores.
          \item Select $x_{t+1}$'s interval.
          \item Encode it, then append $x_{t+1}$.
        \end{enumerate}
      \end{block}
      \vspace{0.4em}
      \centering
      \textcolor{mDarkTeal}{O1 teacher forcing: applicable}\par
      \textcolor{mDarkTeal}{O3 target record: applicable}
    \end{column}
    \hfill
    \begin{column}{0.47\textwidth}
      \begin{block}{Decompression: $x_{t+1}$ is unknown}
        \begin{enumerate}
          \item Recreate scores from decoded prefix.
          \item Find the interval containing code value.
          \item Discover and append $x_{t+1}$.
        \end{enumerate}
      \end{block}
      \vspace{0.4em}
      \centering
      \textcolor{red!70!black}{cannot force an unknown token}\par
      \textcolor{red!70!black}{cannot receive only its interval}
    \end{column}
  \end{columns}
    \vspace{0.9em}
  \begin{block}{Shared invariant}
    The tokenizer, model, context policy, global mask, softmax, and integer quantization must reproduce identical intervals. O2's independent row coders preserve this invariant on both sides.
  \end{block}
\end{frame}



% \begin{frame}{Backup: Prior Work Optimizes Different Objectives}
%   \begin{columns}[T,onlytextwidth]
%     \begin{column}{0.47\textwidth}
%       \begin{block}{Compression research}
%         \begin{itemize}
%           \item improves probability models and ratios
%           \item combines predictors with entropy coding
%           \item rarely reports operator throughput
%         \end{itemize}
%       \end{block}
%       \scriptsize \parencite{schmidhuberSequentialNeuralText1996,goyalDeepZipLosslessData2019,deletangLanguageModelingCompression2024}
%     \end{column}
%     \hfill
%     \begin{column}{0.47\textwidth}
%       \begin{block}{LLM-serving research}
%         \begin{itemize}
%           \item maximizes generated-token throughput
%           \item schedules independent requests
%           \item normally returns sampled tokens, not row-stable intervals
%         \end{itemize}
%       \end{block}
%       \scriptsize \parencite{kwonEfficientMemoryManagement2023,zhengSGLangEfficientExecution2024}
%     \end{column}
%   \end{columns}
% \end{frame}

% \begin{frame}{Backup: Fixed and Variable-Length Codes}
%   \centering
%   \begin{tikzpicture}[
%     card/.style={draw,rounded corners=3pt,minimum width=4.7cm,minimum height=3.0cm,align=center},
%   ]
%     \node[card,fill=blue!7] at (0,0)
%       {\Large\bfseries ASCII\\[0.25em]\fontsize{23}{25}\selectfont\bfseries 8.0 bits\\[-0.05em]per character\\[0.35em]
%        \scriptsize every character gets the same budget};
%     \node[card,fill=mDarkTeal!8] at (5.5,0)
%       {\Large\bfseries Huffman\\[0.25em]\fontsize{23}{25}\selectfont\bfseries $\sim$4.2 bits\\[-0.05em]per character\\[0.35em]
%        \scriptsize frequent characters get shorter codes};
%   \end{tikzpicture}\par
%   \vspace{0.8em}
%   \small Arithmetic coding extends the idea to changing probabilities.\\[-0.1em]
%   \tiny Huffman (1952); Rissanen and Langdon (1979).
% \end{frame}

% \begin{frame}{Backup: Arithmetic Coding Narrows an Interval}
%   \centering
%   \begin{tikzpicture}[font=\sffamily]
%     \node[font=\small\bfseries] at (0,2.5) {start $[0,1)$};
%     \fill[blue!18] (1.5,2.1) rectangle (7.5,2.9);
%     \fill[orange!20] (7.5,2.1) rectangle (10.5,2.9);
%     \fill[red!16] (10.5,2.1) rectangle (11.5,2.9);
%     \draw[thick] (1.5,2.1) rectangle (11.5,2.9);
%     \draw[very thick,mDarkTeal] (1.5,2.02) rectangle (7.5,2.98);
%     \node at (4.5,2.5) {A: .60}; \node at (9,2.5) {B: .30}; \node[font=\scriptsize] at (11,2.5) {C: .10};
%     \node[font=\scriptsize] at (4.5,1.55) {observe A $\Rightarrow [0,.60)$};
%     \draw[-{Stealth[length=6pt]},thick] (4.5,1.35) -- (4.5,0.8);
%     \node[draw,rounded corners=2pt,fill=mDarkTeal!9,minimum width=5.8cm,minimum height=0.85cm] at (4.5,0.25)
%       {repeat inside the retained interval};
%     \node[draw,rounded corners=2pt,fill=mLightBrown!12,minimum width=4.1cm,minimum height=1.35cm,align=center] at (9.5,0.25)
%       {final interval contains\\the compressed binary fraction};
%   \end{tikzpicture}\par
%   \vspace{0.5em}
%   \small Likely symbols retain wider intervals and therefore need fewer bits.
% \end{frame}

% \begin{frame}{Backup: Next-Token Prediction Supplies the Distribution}
%   \centering
%   \Large
%   \[
%     P(x_{1:n})=\prod_{t=1}^{n}P(x_t\mid x_{<t})
%   \]
%   \vspace{0.5em}
%   \[
%     \mathcal{L}_{\mathrm{NTP}}=-\sum_{t=1}^{n}\log P_\theta(x_t\mid x_{<t})
%   \]
%   \vspace{0.6em}

\begin{frame}{Backup: System Under Study}
  \centering
  \includegraphics[width=0.96\textwidth,height=0.78\textheight,keepaspectratio]{figures/tikz/engine_logit_paths.pdf}

  \vspace{0.2em}
  \small Same compressor loop and predictor interface; engine differences enter through score or interval capture.
\end{frame}

% \begin{frame}{Backup: Engine Differences Enter at Score Capture}
%   \centering
%   \includegraphics[width=0.95\textwidth,height=0.75\textheight,keepaspectratio]{figures/tikz/engine_logit_paths.pdf}\par
%   \vspace{0.25em}
%   \small The compressor loop stays fixed; only the engine boundary changes.
% \end{frame}

\begin{frame}{Backup: Each Batch Row Owns a Coder Stream}
  \centering
  \includegraphics[width=0.94\textwidth,height=0.73\textheight,keepaspectratio]{figures/tikz/multistream_pipeline.pdf}\par
  \vspace{0.25em}
  \small Rows run in parallel, but token dependencies remain sequential within each row.
\end{frame}

\begin{frame}{Backup: Full Optimization Ablation}
  \scriptsize
  \renewcommand{\arraystretch}{1.08}
  \begin{tabularx}{\textwidth}{llYYY}
    \toprule
    \textbf{Engine} & \textbf{Configuration} & \textbf{MB/s} & \textbf{Speedup} & \textbf{Coder time} \\
    \midrule
    Transformers & Baseline AC & 0.0018 & $1.00\times$ & 97.1\% \\
    Transformers & Target-interval capture & \metric{0.0347} & \metric{$18.96\times$} & \metric{1.0\%} \\
    Transformers & Multistream + target intervals & 0.0345 & $18.84\times$ & 1.0\% \\
    \midrule
    vLLM & Baseline AC & 0.0031 & $1.00\times$ & 89.9\% \\
    vLLM & Teacher-forced windows & \metric{0.0116} & \metric{$3.43\times$} & 9.2\% \\
    vLLM & Target-interval capture & 0.0065 & $1.94\times$ & 3.0\% \\
    \midrule
    TensorRT-LLM & Baseline AC & 0.0020 & $1.00\times$ & 58.2\% \\
    TensorRT-LLM & Target-interval capture & \metric{0.0044} & \metric{$2.18\times$} & 0.7\% \\
    \bottomrule
  \end{tabularx}
  \vspace{0.45em}
  \tiny Qwen2.5-0.5B, batch 256, context 128, A100, \texttt{text8}; speedup is normalized within each engine.
\end{frame}

\begin{frame}{Backup: Longer Context Mainly Improves Ratio}
  \centering
  \begin{tikzpicture}
    \begin{axis}[
      width=0.88\textwidth,height=0.67\textheight,
      xlabel={Context length},ylabel={Compression ratio},
      xmin=16,xmax=530,ymin=6.1,ymax=7.2,
      xtick={32,64,128,256,512},grid=major,grid style={draw=black!13},
      tick label style={font=\scriptsize},label style={font=\small},
    ]
      \addplot[very thick,mDarkTeal,mark=*] coordinates {(32,6.2288) (64,6.6592) (128,6.9239) (256,7.0600) (512,7.1007)};
    \end{axis}
  \end{tikzpicture}\par
  \vspace{-0.15em}
  \small Throughput remains between 0.0324 and 0.0349 MB/s; ratio gains diminish after the first few hundred tokens.
\end{frame}

\begin{frame}{Backup: Encoding Choice Does Not Fix Inference Cost}
  \small
  \begin{tabularx}{\textwidth}{lYYY}
    \toprule
    \textbf{Encoding} & \textbf{MB/s} & \textbf{Ratio} & \textbf{Interpretation} \\
    \midrule
    Target-interval AC & 0.0349 & \metric{7.1007} & best ratio \\
    Bitpacked ranks & 0.0363 & 2.8735 & ratio collapses \\
    Huffman ranks & 0.0359 & 6.1962 & no speed advantage \\
    \midrule
    LZ4 & \metric{666.66} & 1.83 & CPU baseline \\
    zstd & \metric{586.70} & 3.06 & CPU baseline \\
    FSST & 294.11 & 2.08 & database baseline \\
    xz -9e & 1.34 & 4.30 & slower, stronger ratio \\
    \bottomrule
  \end{tabularx}
  \vspace{0.4em}
  \tiny Cross-system encode comparison; CPU and LLM rows use different hardware and sample sizes.
\end{frame}

\begin{frame}{Backup: Code Compresses Especially Well}
  \centering
  \begin{tikzpicture}
    \begin{axis}[
      ybar,bar width=24pt,width=0.82\textwidth,height=0.62\textheight,
      ymin=0,ymax=12.5,ylabel={Compression ratio},
      symbolic x coords={text8,Python code,Yelp},xtick=data,
      grid=major,grid style={draw=black!13},nodes near coords,
      tick label style={font=\small},label style={font=\small},
    ]
      \addplot[fill=mDarkTeal!75,draw=mDarkTeal] coordinates {(text8,7.739) (Python code,11.203) (Yelp,7.028)};
    \end{axis}
  \end{tikzpicture}\par
  \vspace{-0.2em}
  \small Domain-structured text such as code, logs, and repeated records is a promising target for specialized models.
\end{frame}

\begin{frame}{Backup: What an Engine API Should Expose}
  \begin{enumerate}
    \item \textbf{Row-stable score or interval capture} without a generation round trip.
    \vspace{0.35em}
    \item \textbf{Stateful stepping} that preserves KV state across known prefixes.
    \vspace{0.35em}
    \item \textbf{Reduced outputs} such as target intervals when the consumer does not need dense logits.
    \vspace{0.35em}
    \item \textbf{Deterministic replay} across encoder and decoder deployments.
  \end{enumerate}
\end{frame}

\begin{frame}{Backup: A100 Round-Trip Benchmark}
  \begin{table}
    \centering
    \small
    \begin{tabular}{lrrr}
      \toprule
      Direction & Wall time (s) & Throughput (KiB/s) & Inference (s) \\
      \midrule
      Encode & 6.041 & 0.892 & 0.381 \\
      Decode & 5.108 & 1.055 & 0.393 \\
      \bottomrule
    \end{tabular}
  \end{table}

  \vspace{0.35em}
  \begin{columns}[T,onlytextwidth]
    \begin{column}{0.47\textwidth}
      \textbf{Completed here}
      \begin{itemize}
        \item With small batch size
        \item warm-up, then median of three runs
        \item byte-for-byte round-trip verification
      \end{itemize}
    \end{column}
    \hfill
    \begin{column}{0.47\textwidth}
      \textbf{Still needed cross-codec}
      \begin{itemize}
        \item FSST, OnPair+, and zstd on one CPU core
        \item identical corpus and sample sizes
        \item warmed, amortized jobs for a cost claim
      \end{itemize}
    \end{column}
  \end{columns}

  \vfill
  \tiny 5,517 B of \texttt{text8}; Qwen2.5-0.5B, Transformers + arithmetic coding, batch 256, context 256. Wall time includes startup; compressed artifact is $3.34\times$ smaller.
\end{frame}

